% Generated draft: requires mathematical and source review.
\documentclass[11pt]{article}
\usepackage[margin=25mm]{geometry}
\usepackage{fontspec}
\setmainfont{Latin Modern Roman}
\usepackage{amsmath,amssymb,mathtools,unicode-math}
\setmathfont{Latin Modern Math}
\usepackage{enumitem,needspace,xurl,hyperref}
\setlength{\parindent}{0pt}
\setlength{\parskip}{5pt}
\setlength{\emergencystretch}{3em}
\newcommand{\N}{\mathbb{N}}
\newcommand{\Z}{\mathbb{Z}}
\newcommand{\Q}{\mathbb{Q}}
\newcommand{\R}{\mathbb{R}}
\newcommand{\C}{\mathbb{C}}
\newcommand{\F}{\mathbb{F}}
\newcommand{\A}{\mathbb{A}}
\newcommand{\PP}{\mathbb{P}}
\newcommand{\E}{\mathbb{E}}
\newcommand{\eps}{\varepsilon}
\newcommand{\dd}{\,\mathrm{d}}
\begin{document}
\section*{UnsolvedMath ID: 801066}
\textbf{Source problem code:} catalogue entry
\subsection{Definitions and mathematical statement}
Let \(\mathbb D=\{z\in\C:|z|<1\}\). For \(a\in\mathbb D\setminus\{0\}\), define the dissipative H\'enon map
\[
H_a:\C^2\longrightarrow\C^2,\qquad H_a(x,y)=(ay,x+y^2).
\]
Here dissipativity means that the modulus of the complex Jacobian determinant is less than one; indeed, \(\det DH_a=-a\).

A rotation number \(\theta\in\R/\Z\) is of bounded type if its representative \(\alpha\in(0,1)\) is irrational and its continued fraction expansion \(\alpha=[0;b_1,b_2,\ldots]\) satisfies \(\sup_{n\geq1}b_n<\infty\). Write \(\operatorname{HR}_{\theta}(F)\) for the source assertion that \(F\) has a Herman ring with rotation number \(\theta\). This notation retains an unresolved source qualification: the supplied record does not define a Herman ring for the maps under consideration, including whether the ring must be invariant under the map itself or may be periodic with rotation defined using a return map.

Fix the map denoted by \(H_2\) in the source and a bounded type rotation number \(\theta\). Assume \(\operatorname{HR}_{\theta}(H_2)\). The question is whether
\[
\forall\varepsilon>0\quad\exists a\in\mathbb D\setminus\{0\}\quad
\bigl(|a-1|<\varepsilon\ \text{and}\ \operatorname{HR}_{\theta}(H_a)\bigr).
\]
The map \(H_2\) is not defined in the supplied record. In particular, it is not identified here with a member or an iterate of the displayed family. Consequently, the hypothesis and the Herman ring predicate require clarification from the original source before this becomes a fully specified mathematical statement.
\subsection{Short English statement}
Assuming the source's map \(H_2\) has a Herman ring with bounded type rotation number \(\theta\), can the same rotation number occur on Herman rings of \(H_a(x,y)=(ay,x+y^2)\) for nonzero parameters inside the unit disk arbitrarily close to \(1\)? The meaning of \(H_2\) and the precise Herman ring convention remain unresolved.
\subsection{Sources}
\begin{itemize}
\item Local UnsolvedMath registry, permanent ID \(801066\), source problem code \texttt{catalogue entry}, title ``Henon Maps: Open Problems --- Question 18.'' The supplied record attributes the problem to Berger and Bedford, with proposed year 2023, and labels its status as open; this status is supplied registry information, not an independently verified literature assessment.
\item Berger and Bedford, \emph{H\'enon Maps: A List of Open Problems} (2023), Question 18. Supplied source URL: \url{https://arxiv.org/abs/2312.03907}. The provenance records access on 2026-10-16 and extraction by \texttt{arxiv\_latex}; these are supplied provenance claims, not evidence of independent source consultation for this draft.
\end{itemize}
\end{document}
